\def\level{Première}  % Classe à droite
\def\course{NSI}
\def\chaptername{Boucles FOR et WHILE} % Titre au centre
\def\docname{LP04}        % Identifiant à gauche

\pagestyle{cours_FP}
\makeheader{} % Affiche le bandeau


\begin{easybox}{Les boucles FOR}{}

	La boucle \verb+FOR+ est une boucle bornée. On sait à l'avance combien de fois on va la faire.

\medskip

	\begin{minipage}[b]{0.35\linewidth}
	\begin{lstlisting}[numbers=none]
# En pseudo code
instruction1
Pour i allant de 0 a 6:
	instruction2
instruction3
	\end{lstlisting}
		\end{minipage}\hfill
		\begin{minipage}[b]{0.63\linewidth}
	\begin{itemize}
		\item instruction1 est effectuée
		\item \verb+i+ est un \verb+itérateur+ qui va aller de 0 à 6, c'est à dire 7 tours
		\item On effectuera donc 7 fois l'instruction2 
		\item Quand on a fini la boucle on effectue l'instruction3
	\end{itemize}
		\end{minipage}	
	\end{easybox}
	
	\begin{easybox}{La boucle FOR en python}
		\begin{minipage}[b]{0.33\linewidth}
	\begin{CodePythontexAlt}[TaillePolice=\large, Largeur=1\linewidth,Centre,PremLigne=1]{}
for i in range(0, 7):
	print(i, end=" ")
# affichage: 
# 0 1 2 3 4 5 6
	
				\end{CodePythontexAlt}
		\end{minipage}\hfill
		\begin{minipage}[b]{0.33\linewidth}
	\begin{CodePythontexAlt}[TaillePolice=\large, Largeur=1\linewidth,Centre,PremLigne=1]{}
for i in range(2, 7):
	print(i, end=" ")
# affichage: 
# 2 3 4 5 6
					
				\end{CodePythontexAlt}
		\end{minipage}\hfill
		 \begin{minipage}[b]{0.33\linewidth}
	\begin{CodePythontexAlt}[TaillePolice=\large, Largeur=1\linewidth,Centre,PremLigne=1]{}
for i in range(5):
	print(i, end=" ")
# affichage: 
# 0 1 2 3 4
	
				\end{CodePythontexAlt}
		\end{minipage}
	\end{easybox}
	
	
	\begin{easybox}{Prototypage de la boucle FOR}
	La boucle FOR s'écrit de manière générale de la façon suivante:
	
	\medskip
	
	\begin{minipage}[b]{0.4\linewidth}
	\begin{CodePythontexAlt}[TaillePolice=\large, Largeur=1\linewidth,Centre,PremLigne=1]{}
		for i in range(a, b):
	\end{CodePythontexAlt}
	\end{minipage}\hfill
	\begin{minipage}[b]{0.58\linewidth}
	La boucle est effectuée $b-a$ fois et la variable \verb+i+ prend pour valeur tous les entiers entre $\fo{a}{b}$
	\end{minipage}
	
	
	\end{easybox}
	
	\begin{exemple}[Algorigramme de la structure FOR]{}
	
	
	\begin{center}
	{\verb+for i in range(2, 9):+}
		
	\end{center}
	
	\tcblower{
	
	\begin{CodePythontexAlt}[TaillePolice=\large, Largeur=1\linewidth,Centre,PremLigne=1]{}
		for i in range(2, 9):
			print(i)
		print('fin')
	\end{CodePythontexAlt}
	
	\medskip
	
	
	
	L'itérateur \verb+i+ prend toutes les valeurs entières entre 2 et $9-1$
	
	\renewcommand{\arraystretch}{1.2}
	\begin{center}
	\begin{tabular}{|c|p{2cm}|p{4cm}|}
		\hline
		\verb+i+  & \verb+i < 9+ & instruction \\ \hline{}
		\verb+2+  & \verb+True+ & \verb+print(2)+ \\ \hline{}
		\verb+3+  & \verb+True+ & \verb+print(3)+ \\ \hline{}
		\verb+4+  & \verb+True+ & \verb+print(4)+ \\ \hline{}
		\verb+5+  & \verb+True+ & \verb+print(5)+ \\ \hline{}
		\verb+6+  & \verb+True+ & \verb+print(6)+ \\ \hline{}
		\verb+7+  & \verb+True+ & \verb+print(2)+ \\ \hline{}
		\verb+8+  & \verb+True+ & \verb+print(2)+ \\ \hline{}
		\verb+9+  & \verb+False+ & \verb+print('fin')+ \\\hline
		
	\end{tabular}	
	\end{center}
	
	}
	\end{exemple}
	
	\newpage
	
	\begin{easybox}{Itération sur une chaine de caractères via les indices}{}
	Une chaine de caractères \verb+str+ est un itérateur, c'est à dire qu'il est possible d'itérer sur ses éléments.
	
	\begin{minipage}[b]{0.48\linewidth}
		\begin{ConsolePythontex}[TaillePolice=\large, Label=false]{}
	mot = 'Hello'
	mot[1]
	mot[3]
			\end{ConsolePythontex}
	\end{minipage}\hfill
	\begin{minipage}[b]{0.48\linewidth}
		\begin{center}
			\begin{tabular}{|c|c|c|c|c|c|c|c|}
			\hline
			\textbf{Indice} & \verb+0+ & \verb+1+ & \verb+2+ & \verb+3+ & \verb+4+ \\
			\hline
			\textbf{Caractère} & \verb+H+ & \verb+e+ & \verb+l+ & \verb+l+ & \verb+o+  \\
			\hline
			\textbf{Indice Négatif} & \verb+-5+ & \verb+-4+ & \verb+-3+ & \verb+-2+ & \verb+-1+ \\
			\hline
			\end{tabular}
		\end{center}
	\end{minipage}
	
	
	
	
	\bigskip
	
	Ainsi si on fait varier un itérateur \verb+i+, il est possible d'obtenir chaque élément de \verb+mot+ en appelant \verb+mot[i]+
	
	\begin{minipage}[b]{0.48\linewidth}
	\medskip
	\begin{CodePythontexAlt}[TaillePolice=\large, Largeur=1\linewidth,Centre,PremLigne=1]{}
	mot = 'Hello'
	for i in range(0, len(mot)):
		print(mot[i])
	print('Fin')
			
	\end{CodePythontexAlt}
	\end{minipage}\hfill
	\begin{minipage}[t]{0.48\linewidth}
		\begin{center}
			\renewcommand{\arraystretch}{1}
			\begin{tabular}{|c|c|c|c|}
				\hline
				\verb+i+  & \verb+i < len(mot)+ & instruction & affichage \\ \hline{}
				\verb+0+  & \verb+True+ & \verb+print(mot[0])+ & \verb+H+ \\ \hline{}
				\verb+1+  & \verb+True+ & \verb+print(mot[1])+ & \verb+e+ \\ \hline{}
				\verb+2+  & \verb+True+ & \verb+print(mot[2])+ & \verb+l+ \\ \hline{}
				\verb+3+  & \verb+True+ & \verb+print(mot[3])+ & \verb+l+ \\ \hline{}
				\verb+4+  & \verb+True+ & \verb+print(mot[4])+ & \verb+o+ \\ \hline{}
				\verb+5+  & \verb+False+ & \verb+print('Fin')+ & \verb+Fin+ \\ \hline
			\end{tabular}	
			\end{center}
	\end{minipage}
	\end{easybox}
	
	\begin{easybox}{Itération sur une chaine de caractères via les éléments}{}
	Une chaine de caractères \verb+str+ est un itérateur, c'est à dire qu'il est possible
	d'itérer directement sur les éléments qui constituent cette chaine de caractère
	
	On peut donc créer un variable nommée \verb+lettre+ par exemple qui prendra pour valeur chacun des caractères de \verb+mot+.
	
	\medskip
	
	\begin{minipage}[b]{0.48\linewidth}
	\begin{CodePythontexAlt}[TaillePolice=\large, Largeur=1\linewidth,Centre,PremLigne=1]{}
	mot = 'Hello'
	for lettre in mot:
		print(lettre)
	print('Fin')
	
		\end{CodePythontexAlt}
	\end{minipage}\hfill
	\begin{minipage}[t]{0.48\linewidth}
		\begin{center}
			\renewcommand{\arraystretch}{1}
			\begin{tabular}{|c|c|c|c|}
				\hline
				lettre  & instruction & affichage \\ \hline{}
				\verb+'H'+  & \verb+print(lettre)+ & \verb+H+ \\ \hline{}
				\verb+'e'+  & \verb+print(lettre)+ & \verb+e+ \\ \hline{}
				\verb+'l'+  & \verb+print(lettre)+ & \verb+l+ \\ \hline{}
				\verb+'l'+  & \verb+print(lettre)+ & \verb+l+ \\ \hline{}
				\verb+'o'+  & \verb+print(lettre)+ & \verb+o+ \\ \hline{}
				\verb+None+  & \verb+print('Fin')+ & \verb+Fin+ \\ \hline
			\end{tabular}	
		\end{center}
	
	
	\end{minipage}
	\begin{itemize}
		\item Il n'y a pas besoin de connaitre la longueur du mot
		\item Il n'y  plus l'utilisation d'indice
		\item \verb+lettre+ prend pour valeur chaque caractère de mot
		\item lorsque \verb+lettre+ a pris pour valeur tous les caractères de \verb+mot+, la boucle s'arrête.
	\end{itemize}								
	\end{easybox}
	
	\begin{minipage}[t]{0.45\linewidth}
	\begin{exercice_cours}{Écrire une fonction permettant d'afficher tous les entiers entre 0 et 6 inclus}{}
	\end{exercice_cours}
	\end{minipage}\hspace{15pt}
	\begin{minipage}[t]{0.45\linewidth}
		\begin{exercice_cours}{Écrire une fonction permettant d'afficher tous les entiers entre 2 et 9 inclus}{}
		\end{exercice_cours}
	\end{minipage}
	
	\vspace{1.5cm}
	
	\begin{minipage}[t]{0.45\linewidth}
		\begin{exercice_cours}{Écrire une fonction permettant d'afficher toutes les lettres du mot \verb+'abracadabra'+ en itérant sur les indices}{}
		\end{exercice_cours}
		\end{minipage}\hspace{15pt}
		\begin{minipage}[t]{0.45\linewidth}
			\begin{exercice_cours}{Écrire une fonction permettant d'afficher toutes les lettres du mot \verb+'abracadabra'+ en itérant sur les caractères}{}
			\end{exercice_cours}
	
		\end{minipage}
	
	
	\newpage
	
	\begin{easybox}{Algorithmes de Recherche dans un itérable}
	
	Puisqu'on peut itérer sur une chaine de caractère en observant chaque caractère, on peut décider de faire quelque chose si la caractère a ue particularité.
	
	Ce algorithme compte le nombre de caractères \verb+'a'+ dans la chaine de caractère \verb+mot+
	
	\medskip
	
	\begin{minipage}[t]{0.55\linewidth}
		\begin{tikzpicture}[node distance=1.5cm]
	
	
			% Nœuds
			\node (initialisation_i) [process] at (0, -2){\verb+i = 0+};
			\node (initialisation_compteur) [process] at (0, 0){\verb+compteur = 0+};
			\node (decision1) [decision] at (0, -5){\verb+i+ $<$ \verb+len(mot)+};
			\node (decision2) [decision] at (2, -9){\verb+mot[i] == 'a'+};
			\node (increment_compteur) [process] at (4, -11){\verb-compteur += 1-};
			\node (increment_i) [process] at (6, -5){\verb-i = i + 1-};
			\node (return) [process] at (0, -12){return(\verb+sans_0+)};
		
			\draw [arrow] (initialisation_compteur.south) -- (initialisation_i.north);
			\draw [arrow] (initialisation_i.south) -- (decision1.north);
			\draw [arrow] (decision1.south) -- node[anchor=east]{False}(return.north);
			\draw [arrow] (decision1.east) -| node[anchor=north, yshift=0.5cm] {True}(decision2.north);
			\draw [arrow] (decision2.east) -| node[anchor=north, yshift=0.5cm] {True}(increment_compteur.north);
		
			\coordinate (point1) at (4, -12);
			\coordinate (point2) at ($(initialisation_i.south)!0.5!(decision1.north)$);
			\draw [arrow] (increment_compteur.south) -- (point1);
			\draw [arrow] (decision2.south) |- node[anchor=east, yshift=1cm]{False}(point1);
			\draw [arrow] (point1) -| (increment_i.south);
			\draw [arrow] (increment_i.north) |- (point2);
			
			\filldraw[fill=gray, fill opacity=0.3, draw=black, rounded corners=5mm] 
		(-2,-7) rectangle (7.6,-1);
		\end{tikzpicture}	
	\end{minipage}
	\begin{minipage}[b]{0.43\linewidth}
	
	  \begin{itemize}
			\item on initialise un compteur à 0
			\item on parcourt \verb+mot+ à l'aide d'une boucle \verb+FOR+
			\item on observe donc chaque caractère \verb+mot[i]+
			\item si le caractère est un \verb+'a'+, alors on incrémente le compteur
			\item sinon on ne fait rien 
			\item dans tous les cas on refait une boucle
		\end{itemize}
	
	
	\begin{CodePythontexAlt}[TaillePolice=\large, Largeur=1\linewidth,Centre,PremLigne=1]{}
	def compte_a(mot):
		compteur = 0
		for i in range(len(mot)):
			if mot[i] == 'a':
				compteur=compteur+1
		return compteur
		
	\end{CodePythontexAlt}
	
	La version suivante permet de compter le nombre de \verb+caractere+ dans \verb+mot+
	
	\begin{CodePythontexAlt}[TaillePolice=\large, Largeur=1\linewidth,Centre,PremLigne=1]{}
		def compte(mot, caractere):
			compteur = 0
			for i in range(len(mot)):
				if mot[i] == caractere:
					compteur=compteur+1
			return compteur
			
		\end{CodePythontexAlt}
	\end{minipage}
	
	
	\end{easybox}
	
	
	
	
	
	\begin{minipage}[t]{0.45\linewidth}
		\begin{exercice_cours}{}{}
			En vous inspirant de du programme précédent, proposer un programme permettant de compter le nombre de \verb+lettre_cherchee+ dans \verb+mot+ en itérant sur les éléments.
	
			\begin{CodePythontexAlt}[TaillePolice=\large, Lignes=false, Largeur=1\linewidth,Centre,PremLigne=1]{}
				def compte_lettre_cherchee(mot):
					
	
	
	
	
	
					
	
	
	
	
	
				#
				\end{CodePythontexAlt}
		\end{exercice_cours}
		\end{minipage}\hspace{15pt}
		\begin{minipage}[t]{0.45\linewidth}
			\begin{exercice_cours}{}{}
				\flushleft{Écrire une fonction \verb+est_present(mot, lettre_cherchee)+ qui renvoie \verb+True+ si \verb+lettre_cherchee+ est dans \verb+mot+ et \verb+False+ sinon}
				\begin{CodePythontexAlt}[TaillePolice=\large, Lignes=false, Largeur=1\linewidth,Centre,PremLigne=1]{}
					def est_present(mot, lettre_cherchee):
						
		
		
		
		
		
							
		
		
		
		
		
					#
					\end{CodePythontexAlt}
			\end{exercice_cours}
		\end{minipage}

		\newpage
		
		\section{La boucle WHILE}

		\begin{easybox}{Les boucles WHILE}
		La boucle \verb+WHILE+ est une boucle non bornée. On ne sait à l'avance combien de fois on va la faire.
	
	\begin{minipage}[b]{0.45\linewidth}
	\begin{lstlisting}[numbers=none]
	# En pseudo code
	instruction1
	tant que condition1 vraie:
		instruction2
	instruction3
	\end{lstlisting}
		\end{minipage}\hfill
		\begin{minipage}[b]{0.53\linewidth}
	\begin{itemize}
		\item instruction1 est effectuée
		\item tant que la condition1 est vraie on effectue l'instruction 2
		\item quand la condition1 est fausse, on sort de la boucle;
		\item on effectue l'instruction3
	\end{itemize}
		\end{minipage}	
	\end{easybox}
	
	\begin{easybox}{La boucle WHILE en python}
		\begin{minipage}[b]{0.48\linewidth}
			\begin{CodePythontexAlt}[TaillePolice=\large, Largeur=1\linewidth,Centre,PremLigne=1]{}
				a = 0
				while a < 5:
					print(a, end=" ")
					a = a + 1
				# affichage: 0 1 2 3 4
				\end{CodePythontexAlt}
				
		\end{minipage}\hfill
		\begin{minipage}[t]{0.48\linewidth}
			\begin{CodePythontexAlt}[TaillePolice=\large, Largeur=1\linewidth,Centre,PremLigne=1]{}
				lettre, mot = 'd', 'abracadabra'
				i = 0
				trouve = False
				while not trouve:
					trouve = (lettre == mot[i])
					i = i + 1
				print(i)
				# affichage 7
				# !!! possibilité de boucle infinie
			\end{CodePythontexAlt}
		\end{minipage}\hfill
		 
	\end{easybox}
	
	
	\begin{easybox}{Prototypage de la boucle WHILE}
	La boucle WHILE s'écrit de manière générale de la façon suivante:
	
	\medskip
	
	\begin{minipage}[b]{0.4\linewidth}
	\begin{CodePythontexAlt}[TaillePolice=\large, Largeur=1\linewidth,Centre,PremLigne=1]{}
		while condition1:
			instruction
	\end{CodePythontexAlt}
	\end{minipage}\hfill
	\begin{minipage}[b]{0.58\linewidth}
	La boucle est effectuée tant que \verb|condition1| est vraie. Cette condition est donc un booléen\end{minipage}
	

	
	\end{easybox}
	
	\begin{easybox}{Algorigramme de la structure WHILE}{}
	
	
	\begin{tikzpicture}[node distance=1.5cm]
		% Nœuds
		\node (instruction1) [process] at (0, 0){\verb|a = 0|};
		\node (decision) [decision] at (0, -3){\verb|a| $<$ \verb|5|};
		\node (instruction2) [process] at (0, -7){\verb|print('fin')|};
		\node (instruction3) [process] at (2, -5){\verb|print(a)|};
		\node (instruction4) [process] at (4, -3){\verb|a = a + 1|};
	
		\coordinate (point1) at ($(instruction1.south)!0.5!(decision.north)$);
		\coordinate (point2) at (3, -6);
	
		\draw [arrow] (instruction1.south) -- (decision.north);
		\draw [arrow] (instruction4) |- (point1);
		\draw [arrow] (decision.east) -| node[anchor=north, yshift=0.5cm] {True}(instruction3.north);
		\draw [arrow] (instruction3) |- (point2);
		\draw [arrow] (point2) -| (instruction4.south);
		\draw [arrow] (decision.south) -- node[anchor=east] {False}(instruction2.north);
	
		%\filldraw[fill=gray, fill opacity=0.3, draw=black, rounded corners=5mm] 
		%(-2,-4) rectangle (6,1);
		%\node (instruction2) [process] at (0, -4.5){je prend le parapluie};
		%\node (instruction3) [process] at (0, -6) {je sors de chez moi};
	
		%\coordinate (midpoint1) at ($(instruction2.south)!0.5!(instruction3.north)$);
		%\coordinate (point1) at (2, -4.5);
		% Flèches
		%\draw [arrow] (instruction1) -- (decision);
		%\draw [arrow] (decision.south) -- node[anchor=west] {True} (instruction2.north);
		%\draw [arrow] (instruction2) -- (instruction3);
		%\draw [arrow] (decision.east) -| node[anchor=east, xshift=1.2cm, yshift=-2cm] {False}(point1);
		%\draw [arrow] (point1) |- (midpoint1);
	
		
	\end{tikzpicture}	
	
	\tcblower{
	
	\begin{CodePythontexAlt}[TaillePolice=\large, Largeur=1\linewidth,Centre,PremLigne=1]{}
		a = 0
		while a < 5:
			print(a, end=" ")
			a = a + 1 
		print('fin')
	\end{CodePythontexAlt}
	
	\medskip
	
	
	
	
	\renewcommand{\arraystretch}{1.2}
	\begin{center}
	\begin{tabular}{|c|p{2cm}|p{4cm}|}
		\hline
		\verb|a|  & \verb|a < 5| & instruction \\ \hline
		\verb|0|  & \verb|True| & \verb|print(0)| \\ \hline
		\verb|1|  & \verb|True| & \verb|print(1)| \\ \hline
		\verb|2|  & \verb|True| & \verb|print(2)| \\ \hline
		\verb|3|  & \verb|True| & \verb|print(3)| \\ \hline
		\verb|4|  & \verb|True| & \verb|print(4)| \\ \hline
		\verb|5|  & \verb|False| & \verb|print('fin')| \\ \hline
	\end{tabular}	
	\end{center}
	
	}
	\end{easybox}

	\newpage

	\begin{exercice_cours}{Boucle While}{}
	\begin{enumerate}
		\item On dispose d'une feuille de papier d'épaisseur 0,1 mm. Combien de fois doit-on la plier au minimum pour que l'épaisseur dépasse la hauteur de la tour Eiffel 324 m. Écrire une fonction \verb|nb_plis()| en Python pour résoudre ce problème.
		
		Modifier ce programme pour passer en argument la hauteur à atteindre.

		\item Aujourd'hui un appartement vaut 100 000€. Sa valeur augmente de 1% chaque année.
		\begin{itemize}
			\item 		Écrire un programme en Python pour connaitre sa valeur au bout de 10 ans.
			\item 	Écrire un programme en Python pour savoir au bout de combien d'années sa valeur aura doublé. 
		\end{itemize}

		\begin{minipage}{0.79\linewidth}
				\item  Inès veut construire une pyramide à base carrée comme sur la photo. La pyramide sur la photo a 7 étages. 
			
				\begin{itemize}
					\item 			Inès a 1000 billes. Combien d'étages au maximum aura sa pyramide? Écrire un programme en Python pour répondre au problème. 
					\item 			Modifier ce programme pour déterminer combien d'étages on obtient avec $n$ billes.

				\end{itemize}

		\end{minipage}\hfill
		\begin{minipage}{0.19\linewidth}
						
\includegraphics[width = 1\linewidth]{pyramide.png}

		\end{minipage}

			\item Le 1er janvier, les parents de Gaspard ont déposé sur son livret d'épargne 1500 euros au taux d'intéret de 1,5\% pour qu'il s'achète un scooter à 1700 euros.

			\begin{itemize}
				\item 			Écrire un programme en Python pour savoir au bout de combien d'années il pourra avoir son scooter
				\item 			Modifier le programme pour que l'utilisateur puisse choisir le nombre d'années et la somme initiale.
				\item 			Modifier le programme pour que l'utilisateur puisse également choisir le taux d'intéret en pourcentage.

			\end{itemize}
		
			\item L'ordinateur tire un nombre entier au hasard entre 0 et 100.
			L'utilisateur doit le trouver et pour cela propose des valeurs.
			L'ordinateur indique pour chaque valeur proposée si la valeur est trop petite, trop grande ou s'il a trouvé.
			
			\begin{itemize}
				\item Écrire un programme en Python pour jouer à ce jeu. En combien de coups est-on sûr de trouver?
				\item Modifier le programme pour qu'il s'arrête si l'utilisateur n'a pas trouvé au bout du nombre de coups de la question 2. 	

			\end{itemize}

			\item En mathématiques, la suite de Fibonacci est une suite d'entiers dans laquelle chaque terme est la somme des deux termes qui le précèdent. Elle commence généralement par les termes 0 et 1 (parfois 1 et 1) et ses premiers termes sont 0, 1, 1, 2, 3, 5, 8, 13, 21, 34, 55, 89, etc.
			
			\begin{itemize}
				\item 			Ecrire un programme qui affiche les termes de la suite de Fibonacci jusqu'à ce qu'un des termes soit supérieur strictement à 1000.
				\item 			Modifier ce programme pour obtenir les termes de la suite jusqu'à ce qu'un des termes dépasse $n$ qui sera passé en paramètre.

			\end{itemize}


			\item L'utilisateur donne un entier positif et le programme annonce combien de fois de suite cet entier est divisible par 2.
			
			\item Écrire une fonction qui fait des tirages successifs  d'entiers 0 ou 1 jusqu'à tirer 3 fois 1 consécutivement et
			renvoie le nombre de tirages qui ont été nécessaires.
	
			\item Ecrire une fonction python qui renvoie l'entier $p$ tel que $2^p \leqslant n \leqslant 2^{p+1}$.
		\end{enumerate}
	\end{exercice_cours}


    %%%%%%%%%%%%%%%%%%%%%%%%%%
\begin{easybox}{Les structures conditionnelles IF}

\begin{minipage}{0.32\linewidth}
	\begin{tikzpicture}[node distance=1.5cm]
		% Nœuds
		\node (instruction1) [process] at (0, 0){Instruction 1};
		\node (decision) [decision] at (0, -2){Condition};
		\node (instruction2) [process] at (0, -4.5){Instruction 2};
		\node (instruction3) [process] at (0, -6) {Instruction 3};
	
		\coordinate (midpoint1) at ($(instruction2.south)!0.5!(instruction3.north)$);
		\coordinate (point1) at (2, -4.5);
		% Flèches
		\draw [arrow] (instruction1) -- (decision);
		\draw [arrow] (decision.south) -- node[anchor=west] {True} (instruction2.north);
		\draw [arrow] (instruction2) -- (instruction3);
		\draw [arrow] (decision.east) -| node[anchor=east, xshift=1.2cm, yshift=-2cm] {False}(point1);
		\draw [arrow] (point1) |- (midpoint1);
	\end{tikzpicture}	
	\end{minipage} \hfill
	\begin{minipage}{0.65\linewidth}
		\begin{minipage}{0.55\linewidth}
	\begin{scratch}
    \blocklook{mettre \ovalvariable{total} à \ovalvariable{10}}
\blockif{si \booloperator{\ovalmove{total} > \ovalnum{0}} alors}
{\blocklook{mettre \ovalvariable{total} à \ovaloperator{\ovalvariable{total} + \ovalnum{1}}}
}
\blocklook{Afficher \ovalvariable{total}}
\end{scratch}		
		\end{minipage}\hfill
		\begin{minipage}{0.43\linewidth}
	\begin{lstlisting}[numbers=none]
# En python
Instruction_1
if condition:
	Instruction_2
Instruction_3
	\end{lstlisting}
    
    \begin{lstlisting}[numbers=none]
# En python
total = 10
if total > 0:
	total = total + 1
print(total)
	\end{lstlisting}
		\end{minipage}
	
	
	\begin{itemize}
		\item l'\verb+Instruction 1+ est effectuée
		\item SI la \verb+Condition+ est \verb+True+ alors l'\verb+Instruction 2+ est effectuée
		\item SINON on va directement à l'\verb+Instruction 3+
		\item Dans tous les cas l'\verb+Instruction 3+ est effectuée
	\end{itemize}
	
	La \verb+Condition+ est un test qui est évalué en python par un booléen.
	
	\end{minipage}

\end{easybox}


%%%%%%%%%%%%%%%%%%%%%%ùùù

\begin{exemple}[La condition IF sans ELSE: la valeur absolue d'un nombre]{}
	%Prenons l'exemple d'une fonction calculant la valeur absolue d'une nombre $x$. La valeur absolue d'un nombre est notée en mathématique $|x|$
	
    \begin{minipage}{0.64\linewidth}
        \begin{itemize}
		\item si $x \geqslant 0$, alors $|x| = x$
		\item si $x < 0$, alors $|x|=-x$
	\end{itemize}
        \begin{CodePythontexAlt}[TaillePolice=\large, Largeur=1\linewidth,Centre,Lignes=true,PremLigne=1]{}
	def valeur_absolue(x):
		''' Renvoie |x| '''
		if x < 0:
		# si x < 0 on échange x par -x	
			x = - x
		return x
	\end{CodePythontexAlt}
    \end{minipage}\hfill{}
    \begin{minipage}{0.33\linewidth}
        \begin{center}
			
		\begin{tikzpicture}[node distance=1.5cm]
			% Nœuds
		\node (decision) [decision] at (0, 0){$x < 0$};
		\node (instruction2) [process] at (0, -2){$x = - x$};
		\node (instruction3) [process] at (0, -4) {return $x$};
	
		\coordinate (midpoint1) at ($(instruction2.south)!0.5!(instruction3.north)$);
		\coordinate (point1) at (2, -2.5);
		% Flèches
		\draw [arrow] (decision.south) -- node[anchor=west] {True} (instruction2.north);
		\draw [arrow] (instruction2) -- (instruction3);
		\draw [arrow] (decision.east) -| node[anchor=east, xshift=1.2cm, yshift=-2cm] {False}(point1);
		\draw [arrow] (point1) |- (midpoint1);
	\end{tikzpicture}
	\end{center}
    \end{minipage}
    
	
	
	
	
	
		
		
		
		
	\end{exemple}
	

    


	\begin{easybox}{Les structures conditionnelles IF -- ELSE}{}
		\begin{minipage}{0.4\linewidth}
			\begin{tikzpicture}[node distance=1.5cm]
				% Nœuds
				\node (instruction1) [process] at (0, 0){instruction 1};
				\node (decision) [decision] at (0, -2){Condition};
				\node (instruction2) [process] at (3.5, -4){instruction 3};
				\node (instruction3) [process] at (0, -4) {instruction 2};
				\node (instruction4) [process] at (0, -6) {instruction 4};
				
				\coordinate (midpoint1) at ($(instruction3.south)!0.5!(instruction4.north)$);
					% Flèches
				\draw [arrow] (instruction1) -- (decision);
				\draw [arrow] (decision.east) -| node[anchor=south] {False} (instruction2.north);
				\draw [arrow] (decision.south) -- node[anchor=east, xshift=-0.5cm, yshift=0.3cm] {True} (instruction3.north);
				\draw [arrow] (instruction3) -- (instruction4);
				\draw [arrow] (instruction2) |- (midpoint1);
		
			\end{tikzpicture}		
		\end{minipage}\hfill
		\begin{minipage}{0.58\linewidth}
			\begin{minipage}{0.48\linewidth}
	\begin{scratch}
\blockifelse{si \booloperator{\ovalmove{x} >= \ovalnum{0}} alors}
{\blocklook{mettre \ovalvariable{x} à {\ovalvariable{x ** 2} }}}
{\blocklook{mettre \ovalvariable{x} à {\ovalvariable{x ** 3} }}}
\end{scratch}	
			\end{minipage}\hfill
			\begin{minipage}{0.48\linewidth}
	\begin{lstlisting}[numbers=none]
# En python
	
instruction_1
	
if condition:
	instruction_2
else:
	instruction_3
instruction_4
	\end{lstlisting}		
			\end{minipage}
				
				
	\begin{itemize}
		\item l'\verb+instruction 1+ est effectuée
		\item SI la \verb+Condition+ est \verb+True+ alors l'\verb+instruction 2+ est effectuée
		\item SINON l'\verb+instruction 3+ est effectuée
		\item Dans tous les cas l'\verb+instruction 4+ est effectuée
	\end{itemize}
		\end{minipage}
			
			
	\end{easybox}
	
	\begin{exemple}[La condition IF -- ELSE]{}
		

		\begin{CodePythontexAlt}[TaillePolice=\large, Largeur=1\linewidth,Centre,Lignes=true,PremLigne=1]{}
		def ma_fonction(x):
			if x >= 0:	
				sortie = x ** 2
			else:
				sortie = x ** 3
			return sortie
		\end{CodePythontexAlt}
		\end{exemple}
		
	\begin{easybox}{Les structures conditionnelles IF -- ELIF -- ELSE}{}
	
	\medskip
	
	Dans le cas où il y a plusieurs possibilités de réponse, on peut utiliser deux structures différentes en python.
	
	On suppose une fonction prenant comme paramètre une variable $x$ de type entier.
	
	\[f(x)=1 \texttt{ si }  x>0 \qquad f(x)=0 \texttt{ si }  x=0 \qquad f(x)=-1 \texttt{ si } x<0\]
	
	\begin{minipage}{0.6\linewidth}
		\begin{tikzpicture}[node distance=1.5cm]
			% Nœuds
			\node (instruction1) [process] at (0, 0){Instruction 1};
			\node (decision1) [decision] at (0, -2.3){Condition1};
			\node (decision2) [decision] at (3.5, -4.) {Condition2};	
			\node (instruction2) [process] at (0, -6.5) {Instruction 2};
			\node (instruction3) [process] at (3.5, -6.5){Instruction 3};
			\node (instruction4) [process] at (7, -6.5){Instruction 4};
			\node (instruction5) [process] at (0, -9){Instruction 5};
		
			\coordinate (midpoint1) at ($(instruction2.south)!0.5!(instruction5.north)$);
		
			\draw [arrow] (instruction1) -- (decision1);
			\draw [arrow] (decision1.east) -| node[anchor=west] {False}(decision2.north);
			\draw [arrow] (decision1.south) -- node[anchor=east] {True}(instruction2.north);
			\draw [arrow] (decision2.south) -- node[anchor=east] {True}(instruction3.north);
			\draw [arrow] (decision2.east) -| node[anchor=west] {False}(instruction4.north);
			\draw [arrow] (instruction2) -- (instruction5);
			\draw [arrow] (instruction3.south) |- (midpoint1);
			\draw [arrow] (instruction4.south) |- (midpoint1);
			
	
		
		\end{tikzpicture}	
	\end{minipage}\hfill
	\begin{minipage}{0.38\linewidth}
	
		
		\begin{CodePythontexAlt}[TaillePolice=\large, Largeur=1\linewidth,Centre,Lignes=true,PremLigne=1]{}
			def ma_fonction(x):
				if x > 0:	
					sortie = 1
				elif x == 0:
					sortie = 0
				else:
					sortie = -1
				return sortie
			\end{CodePythontexAlt}
	
			\begin{CodePythontexAlt}[TaillePolice=\large, Largeur=1\linewidth,Centre,Lignes=true,PremLigne=1]{}
				def ma_fonction(x):
					if x > 0:	
						sortie = 1
					else:
						if x == 0:
							sortie = 0
						else:
							sortie = -1
					return sortie
				\end{CodePythontexAlt}
	\end{minipage}
	
	
	\end{easybox}
	
	\begin{exercice_cours}[Compléter les fonctions suivantes]{}
	
        \begin{minipage}{0.48\linewidth}
            \begin{CodePythontexAlt}[TaillePolice=\large, Largeur=1\linewidth,Centre,Lignes=false,PremLigne=1]{}
def plus_grand2(a, b):
    ''' Renvoie le plus grand des 
    deux nombres a et b '''
	# A compléter








    #
	\end{CodePythontexAlt}
        \end{minipage}\hfill{}
        \begin{minipage}{0.48\linewidth}
            \begin{CodePythontexAlt}[TaillePolice=\large,Largeur=1\linewidth,Centre,Lignes=false,PremLigne=1]{}
def plus_grand3(a, b, c):
    ''' Renvoie le plus grand des 
    trois nombres a, b et c '''
    # A compléter
            







    #
		\end{CodePythontexAlt}
        \end{minipage}
        
	
	
	
		
	
	
	\end{exercice_cours}	

\end{document}

\section{L'environnement de travail}

Nous avons à notre disposition plusieurs outils pour écrire du code Python.

\subsection*{Edupython}

\begin{center}
	\includegraphics[width = 0.5\linewidth]{LP01-02.png}
\end{center}

\texttt{Edupython} étant installé localement, il se peut que des obglets soient déjà ouverts. Fermez tous les onglets ouverts et ouvrez un nouvel onglet pour commencer à travailler.

\subsection*{Notebook Capytale}

\begin{minipage}{0.48\linewidth}
	Vous accédez à \texttt{Capytale -- Notebook} depuis votre ENT en cliquant sur \verb|Pédagogie/Capytale|. Le professeur vous donnera un code à entrer dans la zone \verb|Accéder à une activité|.

\end{minipage}\hfill{}
\begin{minipage}{0.48\linewidth}
\begin{center}
	\includegraphics[width = 0.7\linewidth]{LP01-03.png}
\end{center}
\end{minipage}



Vous pourrez également y accéder directement via un lien que vous donnera votre enseignant.

\subsection*{Console Capytale}

\begin{minipage}{0.48\linewidth}
Vous accédez à \texttt{Capytale -- Console} depuis votre ENT de la même manière que pour les Notebook.

\end{minipage}\hfill{}
\begin{minipage}{0.48\linewidth}
\begin{center}
	\includegraphics[width = 0.9\linewidth]{LP01-04.png}
\end{center}
\end{minipage}

\medskip

\begin{center}
\fbox{Quelque soit la façon dont vous tapez du code, n'oubliez pas de sauvegarder TRES régulièrement.}
\end{center}


\pagebreak

\section{Les variables en Python}

\begin{easybox}{Déclarer une variable en Python}

	Pour déclarer une variable en Python, il suffit de lui donner un nom et de lui affecter une valeur à l'aide du symbole $=$. Par exemple, pour déclarer une variable \texttt{x} qui contient la valeur 10, on écrit:

\begin{minipage}{0.48\linewidth}
	\begin{CodePythontexAlt}[TaillePolice=\large, Largeur=1\linewidth,Centre,Lignes=true,PremLigne=1]{}
    a = 10
    \end{CodePythontexAlt}
\end{minipage}\hfill{}
\begin{minipage}{0.48\linewidth}
\begin{scratch}
\blocklook{Affecter à \selectmenu{x} la valeur \ovalnum{10}}\end{scratch}
\end{minipage}
\end{easybox}


\begin{easybox}{Le nom des variables en Python}
	\begin{itemize}
		\item Elles doivent commencer par une lettre ou un underscore (\texttt{\_}).
		\item Elles peuvent contenir des lettres, des chiffres et des underscores.
		\item Elles ne peuvent pas être des mots réservés du langage Python (comme \texttt{if}, \texttt{for}, \texttt{while}, etc.).
		\item \texttt{Variable} et \texttt{variable} sont considérées comme deux variables différentes.
	\end{itemize}
\end{easybox}

\begin{exemple}[Noms de variables valides et invalides]{}

	  \renewcommand{\arraystretch}{1.}
    \begin{tabular}{|l|p{12cm}|}
    \hline
    Bonne variable & \verb+ma_variable+ \\ \hline
    Mauvaise variable & \verb+Ma_variable+,  \verb+ma variable+, \verb+1_variable+, \verb+ variable@+ \\ \hline
    A éviter    & \verb+Variable non explicites+ \\ \hline
    \end{tabular}
\end{exemple}

\begin{easybox}{Afficher une variable en Python}

	Pour afficher la valeur d'une variable en Python, on utilise la fonction \texttt{print()}. Par exemple, pour afficher la valeur de la variable \texttt{x}, on écrit:

\begin{minipage}{0.48\linewidth}
	\begin{CodePythontexAlt}[TaillePolice=\large, Largeur=1\linewidth,Centre,Lignes=true,PremLigne=1]{}
    print(x)
    \end{CodePythontexAlt}
\end{minipage}\hfill{}
\begin{minipage}{0.48\linewidth}
\begin{scratch}
\blocklook{Afficher la valeur de \selectmenu{x}}\end{scratch}
\end{minipage}
\end{easybox}

\begin{easybox}{Les types de variables en Python}
	\begin{itemize}
		\item \textbf{int}: pour les nombres entiers (ex: 1, 42, -5)
		\item \textbf{float}: pour les nombres à virgule flottante (ex: 3.14, -0.001)
		\item \textbf{str}: pour les chaînes de caractères (ex: "Bonjour", 'Python')
		\item \textbf{bool}: pour les valeurs booléennes (ex: True, False)
	\end{itemize}
\end{easybox}

\begin{minipage}{0.19\linewidth}
\begin{lstlisting}[numbers=none]
# type (int)
a = 5
b = 6
\end{lstlisting}
    \end{minipage}\hfill
    \begin{minipage}{0.19\linewidth}    
\begin{lstlisting}[numbers=none]   
# type (str)
a = 'Bonjour'
b = '6'   
\end{lstlisting}    
    \end{minipage}\hfill
    \begin{minipage}{0.19\linewidth}    
\begin{lstlisting}[numbers=none] 
# type (float)
a = 5.0 
b = 6.8
\end{lstlisting}    
    \end{minipage}\hfill{}
\begin{minipage}{0.19\linewidth}    
\begin{lstlisting}[numbers=none]  
# type (bool)
a = True
b = False
\end{lstlisting}    
    \end{minipage}

\begin{easybox}{Opérations mathématiques de base en Python}{}
	\begin{minipage}{0.43\linewidth}
		\begin{itemize}[]
        \item[] \texttt{+} \quad  addition:  \verb~5 + 3~ donne 8
        \item[] \texttt{-} \quad  soustraction \verb~5 - 3~ donne 2
        \item[] \texttt{*} \quad  multiplication \verb~5 * 3~ donne 15
        \item[] \texttt{/} \quad  division \verb~6 / 2~ donne 3.0 (\texttt{float})
\end{itemize}
	\end{minipage}\hfill{}
	\begin{minipage}{0.55\linewidth}
\begin{itemize}[]
        \item[] \texttt{**} \quad  puissance \verb~2 ** 3~ donne 8
        \item[] \texttt{\%} \quad  reste de la division entière \texttt{7 \% 3} donne 1
        \item[] \texttt{//} \quad  quotient de la division entière \texttt{7 // 3} donne 2
    \end{itemize}
	\end{minipage}
\end{easybox}

\begin{minipage}{0.3\linewidth}
    \begin{ConsolePythontex}[TaillePolice=\large, Label=false]{}
        print(17+2)
        print(17-2)
    \end{ConsolePythontex}
\end{minipage}\hfill
\begin{minipage}{0.3\linewidth}
    \begin{ConsolePythontex}[TaillePolice=\large, Label=false]{}
		print(17*2)
		print(17/2)
    \end{ConsolePythontex}
\end{minipage}\hfill
\begin{minipage}{0.3\linewidth}
    \begin{ConsolePythontex}[TaillePolice=\large, Label=false]{}
        print(17%2)
        print(17//2)
    \end{ConsolePythontex}
\end{minipage}

\begin{exercice_cours}[Repérer les les types de variables et leurs valeurs dans les programmes suivants]
	\begin{minipage}{0.48\linewidth}
		\begin{itemize}
		\item \verb~x = 5~
		\item \verb~y = "Hello"~
		\item \verb~z = True~
		\item \verb~a = 3.14~
	\end{itemize}
	\end{minipage}\hfill{}
	\begin{minipage}{0.48\linewidth}
	\begin{itemize}
		\item \verb~c = 6/2~
		\item \verb~d = 6//2~
		\item \verb~f = '1' + '2'~
		\item \verb~g = 2**3~
	\end{itemize}
\end{minipage}
	
\end{exercice_cours}


\medskip

\section{Les variables de type chaine de caractère \texttt{str}}


\begin{easybox}{Les chaines de caractères}{}
    \smallskip
    Les chaines de caractères sont des variables qui contiennent du texte. Elles sont délimitées par des guillemets simples ou doubles.

    \begin{minipage}{0.32\linewidth}
        \begin{ConsolePythontex}[TaillePolice=\large, Label=false]{}
            a = 'Bonjour'
            b = ' a tous'
            c = a + b
            print(c)
        \end{ConsolePythontex}
    \end{minipage}\hfill
    \begin{minipage}{0.32\linewidth}
        \begin{ConsolePythontex}[TaillePolice=\large, Label=false]{}
            a = 'Bonjour'
            b = ' a tous'
            c = a * 3
            print(c)
        \end{ConsolePythontex}
    \end{minipage}\hfill
    \begin{minipage}{0.32\linewidth}
        \begin{ConsolePythontex}[TaillePolice=\large, Label=false]{}
            a = 'Bonjour'
            b = ' a tous'
            c = a[0]
            print(c)
        \end{ConsolePythontex}
    \end{minipage}

    \medskip

    Dans l'exemple ci-dessus, la variable \verb+c+ contient la concaténation de \verb+a+ et \verb+b+ dans le premier cas, la répétition de \verb+a+ trois fois dans le deuxième cas et le premier caractère de \verb+a+ dans le troisième cas.
\end{easybox}

\begin{easybox}{indexation des chaines de caractères}{}
    
    \begin{minipage}{0.5\linewidth}
 Les chaines de caractères peuvent être vues comme un ensemble de caractères ``collés`` les uns aux autres

    Chaque caractère d'une chaine de caractères est indexé: 
    il y a un index qui donne sa position dans la chaine       
    \end{minipage}
    \begin{minipage}{0.5\linewidth}
 \begin{center}
        \begin{tabular}{|c|c|c|c|c|c|c|c|}
        \hline
        \textbf{Indice} & 0 & 1 & 2 & 3 & 4 & 5 & 6 \\
        \hline
        \textbf{Caractère} & B & o & n & j & o & u & r \\
        \hline
        \textbf{Indice Négatif} & -7 & -6 & -5 & -4 & -3 & -2 & -1 \\
        \hline
        \end{tabular}
    \end{center}       
    \end{minipage}
    
   
    

    \begin{minipage}[t]{0.48\linewidth}
        \begin{ConsolePythontex}[TaillePolice=\large, Label=false]{}
            mot = 'Bonjour'
            mot[0]
            mot[4]
        \end{ConsolePythontex}
    \end{minipage}\hfill
    \begin{minipage}[t]{0.48\linewidth}
        \begin{ConsolePythontex}[TaillePolice=\large, Label=false]{}
mot = 'Bonjour'
n = len(mot)
n
mot[-1]
		\end{ConsolePythontex}
    \end{minipage}
   
\end{easybox}

\begin{easybox}{Les opérations sur les chaines de caractères}{}   
    \smallskip
    \begin{itemize}
        \item[] \texttt{+} \quad pour la concaténation:  \verb~'Bonjour' + ' a tous'~ donne \verb~'Bonjour a tous'~
        \item[] \texttt{*} \quad pour la répétition:  \verb~'Bonjour' * 3~ donne \verb~'BonjourBonjourBonjour'~
        \item[] \texttt{[]} \quad pour l'indexation:  \verb~'Bonjour'[0]~ donne \verb~'B'~
        \item[] \texttt{in} \quad pour le test d'appartenance:  \verb~'on' in 'Bonjour'~ donne \verb~True~
        \item[] \texttt{not in} \quad pour le test d'appartenance:  \verb~'on' not in 'Bonjour'~ donne \verb~False~
        \item[] \texttt{len()} \quad pour la longueur:  \verb~len('Bonjour')~ donne \verb~7~
       
    \end{itemize}
\end{easybox}


\begin{minipage}[t]{0.32\linewidth}
    \begin{ConsolePythontex}[TaillePolice=\large, Label=false]{}
        mot1 = 'Bonjour'
        mot2 = 'Thomas'
        mot1 + ' ' + mot2
    \end{ConsolePythontex}
\end{minipage}\hfill
\begin{minipage}[t]{0.32\linewidth}
    \begin{ConsolePythontex}[TaillePolice=\large, Label=false]{}
        'j' in 'Bonjour'
        'b' in 'Bonjour'
    \end{ConsolePythontex}
\end{minipage}\hfill{}
\begin{minipage}[t]{0.32\linewidth}
    \begin{ConsolePythontex}[TaillePolice=\large, Label=false]{}
        'abracadabra'[4]
        len('abracadabra')
    \end{ConsolePythontex}
\end{minipage}

\begin{easybox}{Opérations avancées sur les chaines de caractères: le Slicing}{}
    \hspace{1cm}
    \renewcommand{\arraystretch}{1.2}
    \begin{tabular}{p{2cm}|p{13cm}}
        \texttt{[a:b]} & conserve les caractères d'indice \texttt{[a:b]}:  \verb~'Bonjour'[1:4]~ donne \verb~'onj'~ \\
        \texttt{[a:]} & conserve les caractères d'indice \texttt{[a:]}:  \verb~'Bonjour'[3:]~ donne \verb~'jour'~ \\
        \texttt{[:b]} & conserve les caractères d'indice \texttt{[:b]}:  \verb~'Bonjour'[:3]~ donne \verb~'Bon'~ \\
        \texttt{[::-1]} & inverse les caractères:  \verb~'Bonjour'[::-1]~ donne \verb~'ruojnoB'~ \\
        \texttt{[1:-1]} & élimine le premier et dernier caractère:  \verb~'Bonjour'[1:-1]~ donne \verb~'onjou'~ \\
    \end{tabular}
\end{easybox}

\begin{exercice_cours}[Prédire ce que vont afficher les codes suivants]{}

    \begin{minipage}[t]{0.32\linewidth}        
\begin{lstlisting}[numbers=none]
>>> mot = 'Numerique'
>>> n = len(mot)
>>> mot[5]
>>> mot[n]
>>> mot[-1]
>>> mot[8]
\end{lstlisting}
    \end{minipage} \hfill
    \begin{minipage}[t]{0.32\linewidth}        
\begin{lstlisting}[numbers=none]
>>> mot = 'Numerique'
>>> m in mot
>>> e not in mot
>>> mot+mot[2]
>>> mot*2
>>> 'abc'+mot
\end{lstlisting}
    \end{minipage}\hfill
\begin{minipage}[t]{0.32\linewidth}        
\begin{lstlisting}[numbers=none] 
>>> mot = 'Numerique'
>>> mot[1:5]
>>> mot[:4]
>>> mot[3:]
>>> mot[1:-1]
\end{lstlisting}
    \end{minipage}
\end{exercice_cours}

\begin{easybox}{Principe du transtypage}{}
    Les types des variables en python sont dynamiques, on peut modifier sa valeur ainsi que son type.
    
    \begin{minipage}[t]{0.48\linewidth}
        \begin{ConsolePythontex}[TaillePolice=\large, Label=false]{}
            a = '5'     # a est de type str
            b = int(a)  # b est de type int
            print(b)
        \end{ConsolePythontex}
    \end{minipage}\hfill
    \begin{minipage}[t]{0.48\linewidth}
        \begin{ConsolePythontex}[TaillePolice=\large, Label=false]{}
            a = 4       # a est de type int
            b = str(a)  # b est de type str
            print(b)
        \end{ConsolePythontex}
    \end{minipage}
    
    \end{easybox}

    \section{Les variables de type booléen \texttt{bool}}

    
\begin{easybox}{Les booléens}{}
    \smallskip
    Les booléens sont des variables qui ne peuvent prendre que deux valeurs: \verb+True+ (vrai) ou \verb+False+ (faux).
    
    \smallskip
    Les booléens peuvent être utilisés dans des opérations logiques.
    
	\medskip

    \begin{minipage}{0.32\linewidth}
        \begin{ConsolePythontex}[TaillePolice=\large, Label=false]{}
a = True
b = False
c = a and b
c
        \end{ConsolePythontex}
    \end{minipage}\hfill
    \begin{minipage}{0.32\linewidth}
        \begin{ConsolePythontex}[TaillePolice=\large, Label=false]{}
a = True
b = False
c = a or b
c
        \end{ConsolePythontex}
    \end{minipage}\hfill
    \begin{minipage}{0.32\linewidth}
        \begin{ConsolePythontex}[TaillePolice=\large, Label=false]{}
a = True
c = not a
c
        \end{ConsolePythontex}
    \end{minipage}    
\end{easybox}

\medskip

\begin{minipage}[t]{6cm}
    \begin{center}Table de vérité du \verb+not+\end{center}
    \begin{longtable}[]{@{}cc@{}}
        \hline Expression & Résultat\tabularnewline\hline\endhead{}
        \texttt{not} True & False\tabularnewline{}
        \texttt{not} False & True\tabularnewline\hline{}
    \end{longtable}
\end{minipage}
\begin{minipage}[t]{6cm}
    \begin{center}Table de vérité du \verb+and+\end{center}
    \begin{longtable}[]{@{}cc@{}}
        \hline Expression & Résultat\tabularnewline\hline\endhead{}
        True \texttt{and} True & True\tabularnewline{}
        True \texttt{and} False & False\tabularnewline{}
        False \texttt{and} True & False\tabularnewline{}
        False \texttt{and} False & False\tabularnewline\hline{}
    \end{longtable}
\end{minipage}
\begin{minipage}[t]{6cm}
    \begin{center}Table de vérité du \verb+or+\end{center}
    \begin{longtable}[]{@{}cc@{}}
        \hline Expression & Résultat\tabularnewline\hline\endhead{}
        True \texttt{or} True & True\tabularnewline{}
        True \texttt{or} False & True\tabularnewline{}
        False \texttt{or} True & True\tabularnewline{}
        False \texttt{or} False & False\tabularnewline\hline{}
    \end{longtable}
\end{minipage}

\pagebreak

\begin{easybox}{Booléen comme résultat de comparaison}{}
    Les booléens sont le résultat d'une comparaison. Les opérateurs de comparaison sont:
    \smallskip
    \begin{itemize}
        \begin{minipage}{0.44\linewidth}
            \item[] \texttt{==} \quad pour le test d'égalité
            \item[] \texttt{!=} \quad pour le test de différence
            \item[] \texttt{<} \quad pour le test d'infériorité
        \end{minipage}
        \begin{minipage}{0.44\linewidth}
            \item[] \texttt{>} \quad pour le test de supériorité
            \item[] \texttt{<=} \quad pour le test d'infériorité ou d'égalité
            \item[] \texttt{>=} \quad pour le test de supériorité ou d'égalité
        \end{minipage}
    \end{itemize}

    \end{easybox}


\begin{easybox}{Booléen comme résultat de comparaison}{}  
    \smallskip
    \begin{minipage}{0.32\linewidth}    
        \begin{ConsolePythontex}[TaillePolice=\large, Label=false]{}
a = 5
b = 6
c = (a == b)
c
\end{ConsolePythontex}
    \end{minipage}\hfill
    \begin{minipage}{0.32\linewidth}
        \begin{ConsolePythontex}[TaillePolice=\large, Label=false]{}
a = 5
b = 6
c = (a != b)
c
        \end{ConsolePythontex}
    \end{minipage}\hfill
    \begin{minipage}{0.32\linewidth}
        \begin{ConsolePythontex}[TaillePolice=\large, Label=false]{}
a = 5
b = 6
c = (a < b)
c
        \end{ConsolePythontex}
    \end{minipage}
    
\end{easybox}


\begin{exercice_cours}{    Quelle est la valeur de la variable \verb+c+ après l'exécution des codes suivants?}{}


    \begin{minipage}[t]{0.4\linewidth}
\begin{lstlisting}[numbers=none]   
>>> a = 3
>>> b = 1
>>> c = a > b
\end{lstlisting}

\begin{lstlisting}[numbers=none]
>>> a = True
>>> b = False
>>> c = (a or b) and (not a)
\end{lstlisting}

    \end{minipage} \hfill
    \begin{minipage}[t]{0.58\linewidth}
\begin{lstlisting}[numbers=none]
>>> a, b = True, False
>>> c = (a and b) or (not a) 
\end{lstlisting}   

\begin{lstlisting}[numbers=none]
>>> a = True
>>> b = False
>>> x, y = 7, 9
>>> c = (b or (x == y)) or (a and (x != y))
\end{lstlisting}  
    \end{minipage}

\end{exercice_cours}
